\chapter{زبان‌های نامنظم و لم تزریق}

ماشین‌های متناهی حافظه محدودی دارند، پس نمی‌توانند هر زبانی را تشخیص دهند. در این فصل یاد می‌گیریم چطور ثابت کنیم یک زبان منظم \textbf{نیست}. ابزار اصلی \term{لم تزریق}{Pumping Lemma} است و ابزار کمکی، ویژگی‌های بستاری که در فصل‌های قبل ثابت کردیم.

\section{محدودیت ماشین‌های متناهی}

زبان $B=\{0^n1^n\st n\ge0\}$ را در نظر بگیرید. ماشینی که $B$ را تشخیص دهد باید تعداد $0$ها را به خاطر بسپارد تا بعد آن را با تعداد $1$ها مقایسه کند. اما تعداد $0$ها می‌تواند هر عددی باشد و ماشین فقط تعداد متناهی حالت دارد. پس به نظر می‌رسد هیچ ماشین متناهی‌ای نمی‌تواند $B$ را تشخیص دهد.

آیا این استدلال شهودی کافی است؟ نه! به این دو زبان روی $\{0,1\}$ نگاه کنید:
\begin{align*}
C&=\{w\st w\ \text{\rl{تعداد برابری $0$ و $1$ دارد}}\},\\
D&=\{w\st w\ \text{\rl{تعداد برابری زیررشته $01$ و $10$ دارد}}\}.
\end{align*}
در هر دو به نظر می‌رسد باید «شمردن» انجام شود. اما $C$ منظم نیست (همین فصل ثابت می‌کنیم)، در حالی که $D$ \textbf{منظم است}! دلیلش این است که در هر رشته، هر جا $0$ها به $1$ تبدیل می‌شوند یک $01$ و هر جا $1$ها به $0$ تبدیل می‌شوند یک $10$ داریم. این دو نوع تغییر یکی در میان اتفاق می‌افتند، پس تعدادشان حداکثر یکی فرق دارد و برابرند اگر و تنها اگر رشته با همان نمادی تمام شود که با آن شروع شده (یا تهی باشد). پس
\[D=L(\eps\cup0\cup1\cup0\Sigma^*0\cup1\Sigma^*1).\]
نتیجه: شهود گاهی گمراه‌کننده است و به یک \textbf{اثبات دقیق} نیاز داریم.

\section{لم تزریق}

لم تزریق می‌گوید همه زبان‌های منظم یک ویژگی خاص دارند:

\begin{note}[ایده اصلی]
اگر رشته‌ای از یک زبان منظم از «طول تزریق» بلندتر باشد، بخشی از آن را می‌توان هر چند بار که بخواهیم تکرار کرد (یا حذف کرد) و رشته حاصل باز هم در زبان می‌ماند.
\end{note}

\begin{example}[شهود]
\lr{DFA} زیر زبان $ab^*a$ را تشخیص می‌دهد ($q_3$ حالت تله است):
\begin{automaton}[node distance=2.2cm]
  \node[state,initial] (q0) {$q_0$};
  \node[state,right=of q0] (q1) {$q_1$};
  \node[state,accepting,right=of q1] (q2) {$q_2$};
  \node[state,below=1.3cm of q1] (q3) {$q_3$};
  \path (q0) edge node {$a$} (q1) (q1) edge[loop above] node {$b$} (q1) (q1) edge node {$a$} (q2)
        (q0) edge node[below left] {$b$} (q3) (q2) edge node {$a,b$} (q3) (q3) edge[loop below] node {$a,b$} (q3);
\end{automaton}
رشته‌های زبان: $aa$، $aba$، $abba$، $abbba$، \dots\ هر رشته به طول دست‌کم ۳ مثل $abbbba$ را می‌توان به شکل $x=a$، $y=b$ و $z=bbba$ شکست و برای هر $i\ge0$ داریم $xy^iz=ab^{i}bbba\in L$. دلیل: $y$ دقیقاً روی حلقه $q_1$ خوانده می‌شود و دور زدن چندباره روی این حلقه (یا دور نزدن) ماشین را در همان حالت نگه می‌دارد.
\end{example}

\begin{theorem}[لم تزریق برای زبان‌های منظم]\label{thm:pumping}
برای هر زبان منظم $A$ عددی مثل $p$ (\textbf{طول تزریق}) وجود دارد به طوری که هر رشته $s\in A$ با $\abs{s}\ge p$ را می‌توان به شکل $s=xyz$ نوشت که
\begin{enumerate}
\item برای هر $i\ge0$: \ $xy^iz\in A$،
\item $\abs{y}>0$، و
\item $\abs{xy}\le p$.
\end{enumerate}
\end{theorem}

یادآوری: $y^0=\eps$، پس $xy^0z=xz$. شرط ۲ می‌گوید $y$ تهی نیست (وگرنه قضیه بدیهی می‌شد) و شرط ۳ می‌گوید $y$ در $p$ نماد اول $s$ قرار دارد.

\begin{proofidea}
فرض کنید \lr{DFA} $M$ زبان $A$ را تشخیص دهد و ۵ حالت داشته باشد. اگر $M$ رشته‌ای به طول ۷ را بخواند، از ۸ حالت عبور می‌کند (حالت شروع و ۷ حالت بعدی). طبق \textbf{اصل لانه کبوتری} در ۶ حالت اول (۶ عدد، ۵ لانه) یک حالت تکراری هست؛ مثلاً $q_3$:
\[
s=\ \ s_1\ \ s_2\ \ s_3\ \ s_4\ \ s_5\ \ s_6\ \ s_7,\qquad
q_1\xrightarrow{s_1}\underline{q_3}\xrightarrow{s_2}q_2\xrightarrow{s_3}q_4\xrightarrow{s_4}q_5\xrightarrow{s_5}\underline{q_3}\xrightarrow{s_6}q_5\xrightarrow{s_7}q_2.
\]
پس $x=s_1$، $y=s_2s_3s_4s_5$ (بخشی که ماشین را از $q_3$ به خود $q_3$ برمی‌گرداند) و $z=s_6s_7$. چون $y$ یک «دور» است، می‌توان آن را هر چند بار تکرار یا حذف کرد و ماشین باز هم در همان حالت پایانی تمام می‌کند.
\begin{ltrfig}
\begin{tikzpicture}[jzaut,node distance=2.4cm]
  \node[state,initial] (a) {$q_1$};
  \node[state,right=of a] (b) {$q_3$};
  \node[state,accepting,right=of b] (c) {$q_2$};
  \path (a) edge node {$x=s_1$} (b) (b) edge node {$z=s_6s_7$} (c)
        (b) edge[loop above,looseness=8,min distance=12mm] node {$y=s_2s_3s_4s_5$} (b);
\end{tikzpicture}
\end{ltrfig}
\end{proofidea}

\begin{proof}
فرض کنید \lr{DFA} $M=(Q,\Sigma,\delta,q_1,F)$ زبان $A$ را تشخیص دهد و $p=\abs{Q}$. فرض کنید $s=s_1s_2\cdots s_n\in A$ و $n\ge p$. حالت‌هایی که $M$ هنگام پردازش $s$ از آن‌ها عبور می‌کند را $r_1,r_2,\dots,r_{n+1}$ بنامید؛ یعنی $r_1=q_1$ و $r_{i+1}=\delta(r_i,s_i)$. چون $s\in A$، داریم $r_{n+1}\in F$.

در $p+1$ حالت اول این دنباله ($r_1,\dots,r_{p+1}$) دست‌کم یک حالت تکراری هست، چون فقط $p$ حالت متفاوت داریم (اصل لانه کبوتری). فرض کنید $r_j$ و $r_l$ اولین دو ظهور یک حالت تکراری باشند ($j<l$). چون $r_l$ در میان $p+1$ حالت اول است، $l\le p+1$. قرار دهید
\[x=s_1\cdots s_{j-1},\qquad y=s_j\cdots s_{l-1},\qquad z=s_l\cdots s_n.\]
\begin{itemize}
\item $x$ ماشین را از $r_1$ به $r_j$، $y$ آن را از $r_j$ به $r_l=r_j$ و $z$ آن را از $r_l$ به $r_{n+1}\in F$ می‌برد. چون $y$ از $r_j$ به خود $r_j$ برمی‌گردد، $xy^iz$ هم برای هر $i\ge0$ ماشین را به $r_{n+1}$ می‌برد. پس $xy^iz\in A$.
\item چون $j\neq l$، $\abs{y}=l-j>0$.
\item چون $l\le p+1$، $\abs{xy}=l-1\le p$.\qedhere
\end{itemize}
\end{proof}

\section{روش استفاده از لم تزریق}

لم تزریق را با برهان خلف به کار می‌بریم:
\begin{enumerate}
\item فرض کنید زبان $B$ منظم است. پس طول تزریقی مثل $p$ دارد.
\item رشته‌ای مثل $s\in B$ با $\abs{s}\ge p$ پیدا کنید (معمولاً بر حسب $p$).
\item نشان دهید \textbf{هر} تقسیم $s=xyz$ که شرط‌های ۲ و ۳ را داشته باشد، یک $i$ دارد که $xy^iz\notin B$.
\item این با لم تزریق تناقض دارد، پس $B$ منظم نیست.
\end{enumerate}

این روش را می‌توان مثل یک \textbf{بازی} دید: حریف $p$ را انتخاب می‌کند، شما $s$ را انتخاب می‌کنید، حریف تقسیم $xyz$ را انتخاب می‌کند و شما $i$ را. اگر همیشه بتوانید $i$ را طوری انتخاب کنید که $xy^iz$ از زبان خارج شود، زبان منظم نیست. (این همان نقیض سورهای فصل ۱ است.) انتخاب $s$ مناسب معمولاً به خلاقیت نیاز دارد.

\begin{warning}
لم تزریق می‌گوید هر زبان منظم این ویژگی را دارد، اما \textbf{عکس آن درست نیست}: زبانی که شرایط لم را دارد ممکن است باز هم منظم نباشد (مثالش را در بخش \ref{sec:pumpable} می‌بینیم). پس لم تزریق فقط برای اثبات \emph{نامنظم بودن} به کار می‌رود، نه منظم بودن.
\end{warning}

\section{مثال‌ها}

\begin{example}[$0^n1^n$]\label{ex:0n1n}
$B=\{0^n1^n\st n\ge0\}$ منظم نیست.

فرض کنید $B$ منظم باشد و $p$ طول تزریق آن. رشته $s=0^p1^p\in B$ را انتخاب کنید. هر تقسیم $s=xyz$ را در نظر بگیرید. سه حالت برای $y$ هست:
\begin{enumerate}
\item $y$ فقط $0$ دارد: $xyyz$ تعداد $0$ بیشتری از $1$ دارد، پس در $B$ نیست.
\item $y$ فقط $1$ دارد: $xyyz$ تعداد $1$ بیشتری دارد، پس در $B$ نیست.
\item $y$ هم $0$ و هم $1$ دارد: $xyyz$ ممکن است تعداد برابری $0$ و $1$ داشته باشد، اما ترتیبشان به هم می‌خورد (یک $1$ قبل از یک $0$ می‌آید)، پس در $B$ نیست.
\end{enumerate}
در همه حالت‌ها $xy^2z\notin B$؛ تناقض. (با شرط $\abs{xy}\le p$ می‌شد حالت‌های ۲ و ۳ را از ابتدا کنار گذاشت، چون $y$ باید در $p$ نماد اول، یعنی فقط بین $0$ها باشد.)
\end{example}

\begin{example}[رشته تکراری]
$F=\{ww\st w\in\{0,1\}^*\}$ منظم نیست.

فرض کنید $F$ منظم باشد با طول تزریق $p$. رشته $s=0^p1\,0^p1\in F$ را انتخاب کنید. چون $\abs{xy}\le p$، هم $x$ و هم $y$ فقط از $0$های اول تشکیل شده‌اند: $x=0^m$ و $y=0^k$ با $m+k\le p$ و $k\ge1$. آنگاه $xy^0z=0^{p-k}1\,0^p1$. فرض کنید این رشته به شکل $uu$ باشد. چون $xz$ دقیقاً دو $1$ دارد و به $1$ ختم می‌شود، هر نیمه باید دقیقاً یک $1$ داشته باشد و به $1$ ختم شود. پس نیمه اول $u=0^{p-k}1$ و نیمه دوم $u=0^p1$ است، که برابر نیستند چون $k\ge1$. پس $xz\notin F$؛ تناقض.

(اسلایدهای درس از $s=0^p1^p0^p1^p$ استفاده کرده‌اند که آن هم کار می‌کند. دقت کنید که $s=0^p0^p$ انتخاب بدی است، چون می‌توان آن را تزریق کرد: $y=00$.)
\end{example}

\begin{example}[پالیندروم‌های زوج]
$L=\{ww^R\st w\in\{a,b\}^*\}$ منظم نیست.

فرض کنید $L$ منظم باشد با طول تزریق $p$. رشته $s=a^pb^pb^pa^p=a^pb^{2p}a^p\in L$ را انتخاب کنید. چون $\abs{xy}\le p$، $y=a^k$ با $1\le k\le p$ و فقط از $a$های اول است. آنگاه
\[xy^2z=a^{p+k}b^{2p}a^p.\]
این رشته در $L$ نیست، چون اگر پالیندروم زوج باشد باید تعداد $a$های ابتدا و انتهایش برابر باشد، اما $p+k\neq p$. تناقض.
\end{example}

\begin{example}[جمع اندیس‌ها]
$L=\{a^nb^lc^{n+l}\st n,l\ge0\}$ منظم نیست.

فرض کنید $L$ منظم باشد با طول تزریق $p$. رشته $s=a^pb^pc^{2p}\in L$ را انتخاب کنید. باز هم $y=a^k$ با $1\le k\le p$. با $i=0$:
\[xz=a^{p-k}b^pc^{2p}.\]
برای اینکه این رشته در $L$ باشد باید $2p=(p-k)+p$، یعنی $k=0$، که ممکن نیست. تناقض.
\end{example}

\begin{example}[فاکتوریل]
$L=\{a^{n!}\st n\ge0\}$ منظم نیست.

فرض کنید $L$ منظم باشد با طول تزریق $p$. می‌توانیم فرض کنیم $p\ge2$ (اگر لم برای $p$ برقرار باشد، برای هر عدد بزرگ‌تر هم برقرار است). رشته $s=a^{p!}\in L$ را انتخاب کنید. پس $y=a^k$ با $1\le k\le p$ و
\[xy^2z=a^{p!+k}.\]
اما $p!<p!+k\le p!+p<p!+p\cdot p!=(p+1)!$. (نامساوی آخر چون $p!\ge2$ و پس $p<p\cdot p!$.) پس $p!+k$ اکیداً بین دو فاکتوریل متوالی $p!$ و $(p+1)!$ است و خودش فاکتوریل نیست. یعنی $xy^2z\notin L$؛ تناقض.
\end{example}

\section{استفاده از ویژگی‌های بستاری}

راه دیگر اثبات نامنظم بودن، استفاده از بستاری است:
\begin{enumerate}
\item فرض کنید زبان داده‌شده منظم است.
\item یک عملیات که زبان‌های منظم تحت آن بسته‌اند (اجتماع، الحاق، ستاره، متمم، اشتراک، معکوس) را، احتمالاً همراه با یک زبان منظم دیگر، روی آن اعمال کنید.
\item نشان دهید نتیجه منظم نیست (مثلاً یک زبان شناخته‌شده نامنظم است). این تناقض است.
\end{enumerate}

\begin{example}[تعداد برابر $0$ و $1$]
$C=\{w\st w\ \text{\rl{تعداد برابری $0$ و $1$ دارد}}\}$ منظم نیست.

\textbf{روش ۱ (لم تزریق):} با $s=0^p1^p\in C$، شرط $\abs{xy}\le p$ نتیجه می‌دهد $y$ فقط $0$ دارد، پس $xyyz$ تعداد $0$ بیشتری دارد و در $C$ نیست.

\textbf{روش ۲ (بستاری):} اگر $C$ منظم بود، $C\cap0^*1^*$ هم منظم بود (اشتراک دو زبان منظم). اما $C\cap0^*1^*=\{0^n1^n\st n\ge0\}$ که در مثال \ref{ex:0n1n} دیدیم منظم نیست. تناقض.
\end{example}

\begin{example}[دو نیمه متفاوت]
$D=\{xy\st x,y\in\{0,1\}^*,\ \abs{x}=\abs{y},\ x\neq y\}$ منظم نیست.

فرض کنید $D$ منظم باشد. پس $\overline{D}$ هم منظم است. رشته‌ای به طول فرد در $D$ نیست، پس $\overline{D}=\{ww\st w\in\{0,1\}^*\}\cup\{u\st \abs{u}\ \text{\rl{فرد است}}\}$. با اشتراک با زبان منظم $(\Sigma\Sigma)^*$ (رشته‌های با طول زوج):
\[\overline{D}\cap(\Sigma\Sigma)^*=\{ww\st w\in\{0,1\}^*\}\]
که منظم نیست. تناقض.
\end{example}

\begin{example}
$G=\{0^m1^n\st m\neq n\}$ منظم نیست: اگر بود، $\overline{G}\cap0^*1^*=\{0^n1^n\}$ هم منظم بود.
به همین ترتیب $H=\{w\in\{0,1\}^*\st \text{\rl{تعداد $0$ و $1$ در $w$ برابر نیست}}\}$ منظم نیست، چون $\overline{H}=C$.
\end{example}

\section{یک زبان تزریق‌پذیر که منظم نیست}\label{sec:pumpable}

\begin{example}
زبان $E=\{a^kb^mc^n\st k,m,n\ge0\ \text{\rl{و اگر $k=1$ آنگاه $m=n$}}\}$ را در نظر بگیرید.

\textbf{$E$ شرایط لم تزریق را دارد.} طول تزریق را $p=2$ بگیرید و $s\in E$ را با $\abs{s}\ge2$ در نظر بگیرید. در همه حالت‌ها $x=\eps$ است:
\begin{itemize}
\item $k=0$ و $m=0$: \ $s=c^n$. با $y=c$: \ $xy^iz=c^{n+i-1}\in E$.
\item $k=0$ و $m>0$: \ $s=b^mc^n$. با $y=b$: \ $xy^iz=b^{m+i-1}c^n\in E$.
\item $k=1$ (پس $m=n$): \ $s=ab^nc^n$. با $y=a$: \ $xy^iz=a^ib^nc^n$. برای $i=1$ همان $s$ است و برای $i\neq1$ تعداد $a$ها $1$ نیست و شرطی لازم نیست. پس در $E$ است.
\item $k=2$: \ $s=aab^mc^n$. با $y=aa$: \ $xy^iz=a^{2i}b^mc^n\in E$، چون $2i$ هرگز برابر ۱ نیست.
\item $k\ge3$: \ با $y=a$: \ $xy^iz=a^{k+i-1}b^mc^n\in E$، چون $k+i-1\ge2$.
\end{itemize}
در همه حالت‌ها $\abs{xy}\le2$ و $\abs{y}>0$. پس هر رشته به طول دست‌کم ۲ تزریق‌پذیر است.

\textbf{اما $E$ منظم نیست.} اگر بود، $F=E\cap ab^*c^*=\{ab^nc^n\st n\ge0\}$ هم منظم بود. $F$ را با لم تزریق بررسی می‌کنیم: با $s=ab^pc^p$ و هر تقسیم با $\abs{xy}\le p$ و $\abs{y}>0$: اگر $a$ در $y$ باشد، $xy^0z$ با $a$ شروع نمی‌شود؛ اگر نباشد، $y$ فقط $b$ دارد و $xy^0z$ تعداد برابری $b$ و $c$ ندارد. پس $F$ منظم نیست و در نتیجه $E$ هم منظم نیست.
\end{example}

\section{ویژگی‌های بستاری زبان‌های نامنظم}

\begin{theorem}
اگر $L$ نامنظم باشد، $\overline{L}$ و $L^R$ هم نامنظم‌اند.
\end{theorem}
\begin{proof}
با برهان خلف. اگر $\overline{L}$ منظم بود، چون زبان‌های منظم تحت متمم بسته‌اند، $\overline{\overline{L}}=L$ هم منظم بود؛ تناقض. به همین ترتیب اگر $L^R$ منظم بود، طبق بستاری تحت معکوس $(L^R)^R=L$ منظم بود؛ تناقض.
\end{proof}

\begin{theorem}
دسته زبان‌های نامنظم تحت اجتماع، اشتراک و ستاره \textbf{بسته نیست}.
\end{theorem}
\begin{proof}
برای هر مورد دو زبان نامنظم پیدا می‌کنیم که نتیجه عملیات روی آن‌ها منظم باشد. فرض کنید $A=\{0^n1^n\st n\ge0\}$ که نامنظم است؛ طبق قضیه قبل $\overline{A}$ هم نامنظم است.
\begin{itemize}
\item \textbf{اجتماع:} $A\cup\overline{A}=\Sigma^*$ منظم است.
\item \textbf{اشتراک:} $A\cap\overline{A}=\emptyset$ منظم است.
\item \textbf{ستاره:} قرار دهید $L=A\cup\{0,1\}$. $L$ نامنظم است: اگر منظم بود، $L-\{0,1\}=A$ هم منظم بود (چون $0,1\notin A$ و $\{0,1\}$ منظم است). اما $L^*\supseteq\{0,1\}^*$، پس $L^*=\Sigma^*$ منظم است.\qedhere
\end{itemize}
\end{proof}

\begin{note}
«بسته نبودن» یعنی نتیجه \emph{همیشه} در دسته نمی‌ماند، نه اینکه هیچ‌وقت نمی‌ماند. مثلاً $A\cup A=A$ نامنظم است.
\end{note}

\begin{summary}
\begin{itemize}
\item لم تزریق: در هر زبان منظم، رشته‌های بلند بخشی دارند ($y$، ناتهی و در $p$ نماد اول) که می‌توان آن را تکرار یا حذف کرد. اثبات: اصل لانه کبوتری روی حالت‌های \lr{DFA}.
\item روش استفاده: برهان خلف؛ $s$ را خوب انتخاب کنید و برای \textbf{هر} تقسیم مجاز یک $i$ بد پیدا کنید.
\item عکس لم تزریق درست نیست؛ زبان تزریق‌پذیر هم می‌تواند نامنظم باشد.
\item بستاری هم ابزار قوی‌ای است: با اشتراک با زبانی مثل $0^*1^*$ مسئله را به یک زبان نامنظم شناخته‌شده برگردانید.
\item زبان‌های نامنظم تحت متمم و معکوس بسته‌اند، اما تحت اجتماع، اشتراک و ستاره نه.
\end{itemize}
\end{summary}
